%% Valise à roulettes qu'on traîne sur le quai, en perspective. Repère : z vertical, x le sens
%% de la marche, y transversal. La valise ① a deux mobilités — elle avance et pivote autour de la
%% verticale pour tourner — alors que sa roulette ② n'en a qu'une. Deux nombres différents sur le
%% même dessin.
\documentclass[tikz,border=4pt]{standalone}
\usepackage{liaisons}
%% _objets-common.tex — outils communs aux figures « objet du quotidien avec son repère ».
%%
%% Ces figures servent des exercices où l'axe ne doit PAS être donné dans l'énoncé : l'élève lit
%% l'orientation du repère sur le dessin. D'où deux exigences :
%%   - le repère porte ses TROIS axes, dans une orientation qui change d'une figure à l'autre
%%     (si le vertical était toujours z, « vertical → z » redeviendrait un réflexe) ;
%%   - tout est en \large : à 390 px de large sur un téléphone, \small est illisible.
%%
%% Le \repere de liaisons.sty ne convient pas ici : deux axes seulement, et en \small.
%%
%% RÈGLE DE COULEUR, et elle est stricte : la couleur dit **mobile ou fixe** par rapport au
%% référentiel de l'exercice, jamais « la pièce dont on parle ». Le bâti — sol, route, mur, meuble,
%% tapis, plateau, support — est en solideB (bleu) ou en gris ; TOUT ce qui peut bouger par rapport
%% à lui est en solideA (rouge), même quand plusieurs pièces bougent. Ce sont les pastilles ① ② qui
%% désignent la pièce interrogée, c'est leur seule raison d'être. Colorer « la pièce ① » en rouge et
%% « la pièce ② » en bleu induirait en erreur une fois sur deux l'élève qui a appris que le rouge
%% bouge — le cas s'est produit sur le vélo, dont le cadre était bleu alors qu'il roule.

%% \repereXYZ[longueur]{point}{angle/label, angle/label, angle/label}
%% Trois flèches depuis un même point, chacune à l'angle voulu. L'axe « qui sort de la feuille »
%% se dessine comme les autres, en diagonale, à la manière des perspectives du cours.
\newcommand\repereXYZ[3][0.95]{%
  \foreach \ang/\lab in {#3} {%
    \draw[-{Stealth[length=6pt]}, line width=0.8pt] #2 -- ++(\ang:#1)
      node[pos=1, anchor={\ang+180}, inner sep=3pt, font=\large] {$\vec{\lab}$};%
  }}

%% \numero{point du repère}{point visé sur le dessin}{n} — pastille noire reliée à la pièce.
\newcommand\numero[3]{%
  \draw[line width=0.5pt, black!55] #1 -- #2;
  \node[circle, fill=black, text=white, inner sep=2pt, font=\large\bfseries] at #1 {#3};}

%% Épaisseurs : le dessin doit tenir à la taille d'un téléphone, donc peu de traits et des traits gras.
\tikzset{
  objet/piece/.style={line width=1.6pt, line cap=round, line join=round, draw=solideA},
  objet/bati/.style={line width=1.6pt, line cap=round, line join=round, draw=solideB},
  objet/fin/.style={line width=0.7pt, line cap=round, draw=black!55},
}

\begin{document}
\begin{tikzpicture}[x=1cm,y=1cm]

% --- le sol : le bâti
\draw[line width=1.6pt, line cap=round, black!55] (-0.8,0) -- (5.6,0);
\foreach \h in {-0.6,-0.1,...,5.3} { \draw[objet/fin] (\h,0) -- ++(-0.3,-0.32); }

% --- la caisse, inclinée comme on la traîne
\begin{scope}[rotate around={-14:(1.9,0.3)}]
  \coordinate (ca) at (1.15,0.3);
  \coordinate (cb) at (2.65,0.3);
  \coordinate (cc) at (2.65,2.5);
  \coordinate (cd) at (1.15,2.5);
  \draw[objet/piece, fill=solideA!7] (ca) -- (cb) -- (cc) -- (cd) -- cycle;
  % la profondeur en déplacements relatifs : dans un scope tourné, un point nommé du dehors
  % serait ajouté sans son orientation
  \draw[objet/piece] (cb) -- ++(30:0.75) coordinate (pb) -- ++(0,2.2) coordinate (pc) -- (cc);
  \draw[objet/piece] (pc) -- ++(-1.5,0);
  \draw[objet/piece] (cd) -- ++(30:0.75);
  % la poignée télescopique
  \draw[objet/piece] ($(cd)+(0.28,0)$) -- ($(cd)+(0.28,1.15)$) -- ($(cc)+(-0.28,1.15)$) -- ($(cc)+(-0.28,0)$);
  % la roulette, au pied
  \coordinate (ro) at ($(ca)+(0.3,-0.18)$);
  \draw[objet/piece, fill=white] (ro) circle (0.28);
  \fill[solideA] (ro) circle (0.07);
\end{scope}

% --- repère : y vertical, x le sens de la marche, z transversal
\repereXYZ{(4.85,0.9)}{90/z, 0/x, 30/y}

\numero{(0.45,2.75)}{(1.35,2.15)}{1}
\numero{(0.55,-0.95)}{(ro)}{2}

\end{tikzpicture}
\end{document}
